\documentclass[a4paper,12pt]{article}

\usepackage[utf8]{inputenc} 
\usepackage{amsmath}
\usepackage[a4paper, margin=0.9in]{geometry} % Adjust the margin here
\usepackage{multirow}
\begin{document}

\section*{Question 1}
Two particles are moving uniformly in two-dimensional space on straight lines. The first particle starts at point $\vec{r}_1 = (0, 4) \, \text{m}$ with velocity $\vec{v}_1 = (1, -2) \, \text{m/s}$, and the second particle starts at point $\vec{r}_2 = (-3, 2) \, \text{m}$ with velocity $\vec{v}_2 = (2, 1) \, \text{m/s}$. Find:
\begin{enumerate}
    \item[(a)] The vector of their relative position with respect to time, $\vec{R}(t)$. --- \textit{0.5 points}
    \item[(b)] The moment $t_c$ of closest approach and its distance. --- \textit{0.3 points}
    \item[(c)] The intersection point of their trajectories. --- \textit{0.2 points}
\end{enumerate}




\section*{Question 2}
An object is thrown vertically downward from a height $h$ with an initial velocity $v_0$. The motion takes place on earth with gravitational acceleration $g$ and air resistance is neglected.
\begin{enumerate}
    \item[(a)] What should be $v_0$ so that the object lands $\tau$ seconds earlier than it would if it was released from rest? --- \textit{0.7 points}
    \item[(b)] Assume that there was also a force acting in the perpendicular direction as the throw of the object and accelerating it as $a=\frac{g}{\sqrt{3}}$. What would be the range of the object in this case? --- \textit{0.3 points}
\end{enumerate}

\section*{Question 3}
A particle of mass $m$ can move without friction along the  $x$-axis. At time $t = 0$, it is at rest at position $x = 0$. A force $F_x(t) = a + bt$ (where $a$ and $b$ are constants) acts on it. 
\begin{itemize}
    \item[(a)] Determine the acceleration $a(t)$, the velocity $v(t)$, the displacement $x(t)$. --- \textit{0.6 points}
    \item[(b)] Determine the power delivered by the force at time $t$. --- \textit{0.4 points}
\end{itemize}

\section*{Question 4}
A simple pendulum with mass $m$ and length $l$, initially displaced by an angle $\phi_0$ from the vertical and given an initial tangential speed $v_0$ at $t = 0$. Find the oscillation amplitude $\phi_1$. --- \textit{1 point}

\end{document}
